% =============================================================================
% SECTION 1: INTRODUCTION
% =============================================================================

\section{Introduction}
\label{sec:intro}

Large language models reason through generation: each token is sampled from a
distribution conditioned on all previous
tokens~\citep{brown2020gpt3,touvron2023llama}. This process produces fluent text, but
fluency is not factual accuracy. LLMs hallucinate, generating confident claims that
contradict source material or cannot be verified against any authoritative
source~\citep{ji2023survey,huang2023survey}. The problem worsens with scale: larger
models hallucinate more convincingly, not less~\citep{lin2022truthfulqa}.

Knowledge graph-constrained decoding solves this by replacing the LLM's free-form
generation with a structurally bounded traversal of verified facts. GCR~\citep{luo2025-graph-constrained-reasoning},
the leading framework, achieves zero hallucination on KGQA benchmarks by encoding all
reasoning paths within $L$ hops of the question entities in a prefix trie and blocking
any token that deviates from this structure. The constraint is absolute. Every generated
token must trace back to a real triple on the graph.

The constraint is also indifferent to what the question asks. A three-hop query with
twenty out-degree relations and three question entities admits roughly 24{,}000
structurally valid paths. Only one leads to the correct answer. The LLM must navigate
among the remaining 23{,}999 using the same internal knowledge that produced
hallucination in the first place. Tightening the constraint, filtering out irrelevant
paths, seems like the obvious fix. We show that it is not.

\subsection{Background}
\label{sec:intro:background}

LLMs generate text autoregressively: each token is sampled from a distribution
conditioned on all previous tokens~\citep{brown2020gpt3,touvron2023llama}. This
process produces fluent text, but fluency is not factual accuracy. LLMs hallucinate,
generating confident claims that contradict source material or cannot be verified
against any authoritative source~\citep{ji2023survey,huang2023survey}. The problem
worsens with scale: larger models hallucinate more convincingly, not less. Because
autoregressive generation has no internal mechanism for fact-checking, an initial error
compounds across subsequent steps, producing persuasive yet false reasoning
chains~\citep{huang2023survey}.

Knowledge Graph Question Answering (KGQA) grounds reasoning in structured facts. A
knowledge graph stores facts as triplets, $(\text{entity}, \text{relation},
\text{entity})$, and KGQA requires traversing chains of these triplets to answer
natural language questions~\citep{bollacker2008freebase}. Benchmarks such as
WebQSP~\citep{yih2016websqp} and Complex WebQuestions~\citep{talmor2018cwq} evaluate
multi-hop path traversal over Freebase, demanding strict adherence to graph topology.

\subsection{The Permissiveness Problem}
\label{sec:intro:gap}

Retrieval-Augmented Generation (RAG) mitigates hallucination by supplying retrieved
evidence as context during inference~\citep{lewis2020rag,asai2024selfrag}. But retrieved
evidence acts as a soft prompt, not an absolute boundary. LLMs can ignore relevant
context or generate unsupported claims beyond the retrieved
information~\citep{asai2024selfrag}.

Graph-constrained decoding enforces a harder boundary. By applying a validity mask over
the output vocabulary at each decoding step, a constraint oracle guarantees that every
generated token corresponds to a valid path on the knowledge graph. GCR operationalizes
this using a pre-computed KG-Trie as the constraint oracle during beam
search~\citep{luo2025-graph-constrained-reasoning}. The trie encodes all reasoning
paths within $L$ hops of the question entities as formatted strings, and a
prefix-lookup function blocks any token sequence that deviates from these strings.
Hallucination in the reasoning path drops to zero.

The trie, however, is static. Constructed once before decoding and never updated, its
constraint set, the set of tokens the LLM is allowed to generate, depends only on
the graph structure and the question entities:
\begin{equation}
  W_{\text{val}}^{\text{GCR}}(t) = f(G, E_q)
  \label{eq:gcr_constraint}
\end{equation}
It does not depend on the question $q$ or the partial reasoning chain $y^{<t}$. For a
three-hop question with average out-degree $d \approx 20$ and three question entities, GCR
admits roughly $3 \times 20^3 = 24{,}000$ structurally valid paths. The LLM must
distinguish the relevant paths from the irrelevant ones using the same parametric
knowledge that produced hallucinations before constrained decoding was
introduced~\citep{luo2025-graph-constrained-reasoning,li2024-decoding-on-graphs}.

\subsection{Existing Approaches and Their Constraints}
\label{sec:intro:existing}

The tight coupling between LLMs and knowledge graphs has been pursued along two threads.
Retrieval-based methods train specialized subgraph retrievers to extract
question-focused subgraphs from the complete
KG~\citep{luo2024rog,mavromatis2022rearev}. The concentrated subgraph is then fed to
the LLM for answer prediction. These methods shift the burden of graph traversal to the
retriever while the LLM merely selects from a pre-filtered set of candidates. Training
the retriever requires annotated data, and the trained retriever may struggle with
out-of-domain scenarios where the KG taxonomy differs from the training
distribution~\citep{li2024-decoding-on-graphs}.

Agent-based methods take the opposite approach: the LLM itself interacts with the KG
iteratively, executing operations such as neighbor exploration and path
selection~\citep{jiang2023structgpt}. These methods demand substantial computational
resources and place high demands on the LLM's instruction-following ability, which
limits their applicability to small-scale models~\citep{li2024-decoding-on-graphs}.

Constrained decoding methods occupy a middle ground. GCR converts the KG into a
trie-based index and constrains the LLM's output to valid paths, achieving zero
hallucination in the reasoning
trace~\citep{luo2025-graph-constrained-reasoning}. DoG~\citep{li2024-decoding-on-graphs}
extends this by dynamically expanding the local trie as reasoning proceeds, adapting the
constraint set to the partially generated chain. Both methods, however, share a common
assumption: every structurally valid path should be admitted. Within $L$ hops, GCR
admits everything. At each step, DoG admits every topologically adjacent path. Neither
filters for semantic relevance: whether a path is actually related to the question
being asked.

\subsection{DCA-Trie: Dynamic Context-Aware Constrained Decoding}
\label{sec:intro:method}

We introduce DCA-Trie, a dynamic constrained decoding framework that conditions the
valid token set on the knowledge graph, the input question, and the partially generated
reasoning chain:
\begin{equation}
  W_{\text{val}}^{\text{DCA}}(t) = f(G, E_q, q, y^{<t})
  \label{eq:dca_constraint}
\end{equation}
The constraint set adapts at each decoding step. A lightweight symbolic module, the
TypeOracle, checks two ontology gates per candidate triple. The range gate verifies
that the tail entity's type matches the relation's range at every hop. The type gate
verifies that the terminal entity's type matches the inferred answer type at the final
hop. Paths failing either gate are pruned before the LLM sees them.

Unlike embedding-based relevance scorers, the TypeOracle requires no neural encoder, no
learned threshold, no per-dataset calibration, and no GPU memory for encoder inference.
Gate evaluation is $O(1)$ per candidate triple, adding negligible overhead to the
decoding process.

\subsection{Results and Contributions}
\label{sec:intro:results}

We introduce the Semantic Irrelevance Ratio (SIR), a metric that measures the fraction
of structurally valid paths that are semantically irrelevant to the question. SIR
separates constraint quality from answer accuracy: a system can have excellent SIR
(tight constraints) and poor Hits@1 (wrong answer), or poor SIR (loose constraints) and
excellent Hits@1 (correct answer). Before SIR, these two dimensions were entangled.

On WebQSP, DCA-Trie reduces SIR from 1.0 to 0.145, an 85.5\% removal of
irrelevant paths, while maintaining a false negative rate below 5\%. The TypeOracle
that drives this filtering is purely symbolic: two ontology gates, $O(1)$ per candidate
triple, no neural encoder, no learned threshold.

SIR also reveals that constraint tightness and answer accuracy are independently
controlled variables. Tighter constraints reduce the search space but do not improve
Hits@1. Four mechanisms govern this relationship: gold-path exclusion (the TypeOracle
occasionally removes the correct path), beam competition (fewer paths reduce diversity
in beam search), type-inference noise (the regex-based question classifier has
imperfect accuracy), and the precision-recall trade-off (tighter constraints improve
precision but reduce recall). The result holds across ablation conditions.

The contributions of this work are:
\begin{enumerate}[label=(\arabic*)]
  \item We introduce the Semantic Irrelevance Ratio (SIR), a metric that separates
        constraint quality from answer accuracy, enabling independent evaluation of what
        constrained decoding filters versus what the LLM actually uses.
  \item We propose DCA-Trie, a dynamic constrained decoding framework that conditions
        the valid token set on the question and partial output via a TypeOracle with two
        ontology gates, reducing SIR by 85.5\% with negligible latency overhead.
  \item We establish that constraint tightness and answer accuracy are independently
        controlled variables, and characterise the four mechanisms (gold-path exclusion,
        beam competition, type-inference noise, and the precision-recall trade-off) that
        govern their relationship.
\end{enumerate}

The remainder of this paper is organized as follows. Section~\ref{sec:related} reviews
related work on hallucination mitigation, knowledge graph reasoning, and constrained
decoding. Section~\ref{sec:method} presents the DCA-Trie framework, including the
TypeOracle, the two-gate mechanism, and the SIR metric.
Section~\ref{sec:theoretical} provides theoretical analysis of correctness and
search-space reduction. Section~\ref{sec:experiment_design} describes the experimental
design, and Section~\ref{sec:results} reports results. Section~\ref{sec:analysis}
presents extended analysis including ablation and sensitivity studies.
Section~\ref{sec:discussion} discusses implications and limitations, and
Section~\ref{sec:conclusion} concludes.
